\scanpage{21}
tel que $df=0$, \inserted{donc $f\in K^{*p}$}, on a
$f\in\underline{\mathcal O}_X^p$.
\inserted{[$\underline{\mathcal O}_X^*/\underline{\mathcal O}_X^{*p}\longrightarrow
\mathcal R_X^*/\mathcal R_X^{*p}$ est injectif]}
Donc $X$ \uline{étant normal}, on a un hom \uline{injectif}
\diagram{\begin{tikzcd}
0 \arrow[r] & \underline{\mathcal O}_X^*/\underline{\mathcal O}_X^{*p} \arrow[r] & Z\underline\Omega^1_{X/k}
\end{tikzcd}}
D'après Cartier, un élément dans l'image est tel que
$C\omega=\omega$ [$C$ l'opération de Cartier], et du pt de
vue birationnel la réciproque est vraie : si $\omega$ est une
1-forme fermée telle que $C\omega=\omega$, alors $\omega=\dfrac{df}{f}$,
où $f$ est une \uline{fonction rationnelle} qui n'est nulle
sur aucun \unsure{composant|composante} de $X$,
\deleted{D'ailleurs, comme}\inserted{déterminée mod} $K^{*p}$.
Comme $\underline{\mathcal O}_X^*/\underline{\mathcal O}_X^{*p}\longrightarrow
\underline{\mathcal R}_X^*/\underline{\mathcal R}_X^{*p}$ est injectif,
pour montrer que $\omega$ provient d'une section de
$\underline{\mathcal O}_X^*/\underline{\mathcal O}_X^{*p}$, la question est \uline{locale},
donc il faut prouver que si $f\in(K=\text{corps des fractions de }\mathcal O)^*$,
telle que $\dfrac{df}{f}$ \uline{régulière}, alors $f=g^p h$, avec $h\in\mathcal O^*$,
i.e. \deleted{un diviseur} l'hom
\[
\mathcal R_X^*/\underline{\mathcal O}_X^*\longrightarrow
\underline\Omega^1_{\unsure{k}}(\underline K)/\underline\Omega^1(\underline{\mathcal O}_X)
\]
est \uline{injectif}. $X$ étant normal, \deleted{donc régulier en \unsure{codim 1}},
on est ramené au cas où $X$ est \deleted{\unsure{non}} régulier ...
\mspage{20}
